\section{Background and Related Work}
This scoped narrative synthesis establishes context for ArchSync's design. Sources were selected purposively from the retained keyword-search corpus for architecture conformance, drift detection, and static-analysis reliability, emphasizing 2022--2025. Four earlier works supply concept origins.

\textbf{Erosion and reconstruction.} A mapping of architecture erosion describes its manifestations and consequences \cite{li2022erosion}. Code-review studies identify symptoms in OpenStack and later OpenStack/Qt projects \cite{li2022symptoms,li2023warnings}. Practitioner evidence identifies documentation, responsibilities, and shared practices as ways to manage drift \cite{anthony2024drifting}. Reconstruction taxonomies organize recovery strategies \cite{ducasse2009reconstruction,sar2024taxonomy}, while a microservice mapping compares tool purposes and availability \cite{bakhtin2024mapping}. View-based analysis compares intended and reconstructed specifications in an industrial case \cite{uzun2024drift}. ArchSync restricts its observation boundary to explicit service and infrastructure relationships.

\textbf{Conformance.} Reflexion Models, static conformance comparisons, and DCL provide the origins of model comparison and allowed/forbidden dependency rules \cite{murphy1995reflexion,knodel2007comparison,terra2009dcl}. CALint applies dependency-rule checks to Python \cite{calint2022}. ArchLintor is a close TypeScript comparison, detecting module-interaction violations and proposing CI/CD integration \cite{archlintor2025}. ArchSync's typed service topology differs from that module boundary. IaC metrics address deployment coupling \cite{ntentos2022iac}, and 3Erefactor searches for executable inconsistency-reducing repairs \cite{refactor2024}. LLM-generated architectural tests distinguish execution success from correct rule expression \cite{lucena2025gatekeepers}. ArchSync executes explicit rules deterministically.

\textbf{Change-time feedback.} Commit-level prediction studies architectural change \cite{commits2025}; repository mining tracks microservice evolution \cite{evolution2024}; industrial conformance visualization supplies practitioner evidence \cite{thermofisher2025}. ArchSync distinguishes introduced violations from architectural evolution at a Git head.

\textbf{Evaluation reliability.} Replication support remains an evaluation concern \cite{konersmann2022replicability}. Decomposition surveys cover an adjacent problem \cite{abgaz2023decomposition}, whereas a microservice-recovery comparison executes nine tools and reports individual and combined recovery performance \cite{schneider2025comparison}. That study supplies an external-comparison precedent. Static visual models and multiple dependency models motivate future multi-source work \cite{staticmodels2024,dependencies2024}. Interrogation testing exposes both precision and soundness failures \cite{kaindlstorfer2024interrogation}, motivating our explicit hard negatives.


\textbf{Comparison with reported external evidence.} Table~\ref{tab:external} separates the closest rule-checking approaches from a broader recovery benchmark. The entries summarize published observations. ArchLintor's Corrector.ts example contains one inserted absence, two divergences and one unused-dependency alert; the alert is not a violation, and no precision/recall benchmark is reported \cite{archlintor2025}. CALint's publisher abstract reports three case studies; its unavailable full text prevents a finer result audit \cite{calint2022}.

\begin{table}[t]
\caption{Published comparison context and our evidence boundary. Different languages, units and datasets prevent numerical ranking.}
\label{tab:external}
\centering\small
\begin{tabular}{@{}p{.23\columnwidth}p{.33\columnwidth}p{.35\columnwidth}@{}}
\toprule
Approach & Unit and decision & Evaluation evidence\\\midrule
ArchLintor \cite{archlintor2025} & TypeScript module dependencies; allowed, forbidden, required & One author-developed eight-module example; inserted situations\\
CALint \cite{calint2022} & Python dependencies; Clean Architecture rule & Three cases reported in abstract; full text unavailable\\
Recovery comparison \cite{schneider2025comparison} & Components, connections, endpoints; recovery correctness & Nine tools on 17 Java/Spring applications\\
ArchSync (this work) & Typed service relations; rules and introduced findings & 20 co-developed patches; 40 detector signals; no independent holdout\\\bottomrule
\end{tabular}
\end{table}

For a concrete external result, Schneider et al.'s Table~6(d) reports Code2DFD TP=584, FP=45 and FN=143, with precision 0.93, recall 0.80 and F1 0.86 \cite{schneider2025comparison}. These scores measure recovered characteristics on Java/Spring applications. ArchLintor provides a closer language-level comparison through explicit module rules and source-linked diagnostics. ArchSync integrates typed service observations with introduced-finding decisions; Section~6 addresses the evaluation boundaries across these approaches.
